\documentclass[aps,prd,twocolumn,superscriptaddress,nofootinbib,10pt]{revtex4-2}

\usepackage{amsmath,amssymb}
\usepackage{graphicx}
\usepackage[colorlinks=true,citecolor=blue,linkcolor=blue,urlcolor=blue]{hyperref}
\usepackage{bm}

\begin{document}

\title{A ratio-locked radio signature of a universal shift in the electron mass}

\author{Justin Pulford}
\email{pulfordj420@gmail.com}
\affiliation{Unaffiliated}

\date{\today}

\begin{abstract}
A universal fractional shift $\varepsilon$ in the electron mass would imprint fixed relative
shifts across radio bands, with ratios set by atomic physics and independent of $\varepsilon$
itself. That multi-band \emph{pattern}---and the falsification it enables---is the content of this
paper, not a competitive bound on the amplitude. Methanol absorption already constrains
$|\Delta\mu/\mu|$ at the few$\times 10^{-7}$ level, some thirty-five times tighter than the best
$21\,$cm reinterpretation quoted below; a single molecular comparison cannot, however, test whether
a shift lies in $m_e$ rather than in a coupling, nor whether every band shares it. We derive how
such a shift would \emph{propagate}: $+2$ for the hydrogen $21\,$cm hyperfine frequency, $+1$ for
radio recombination lines, $-1$ for a dispersion-measure reconstruction (degenerate with the fitted
column unless an independent electron column is supplied), $-1$ or $-3$ for synchrotron
characteristic frequency according as the emitting population is labelled by Lorentz factor or by
energy, and $-2$ for Faraday rotation. Only the $21\,$cm and Faraday rows are presently usable in a
Gaussian multi-band fit; the other three are theoretical propagation channels, not simultaneous
precision measurements available today. The pattern separates a varying electron mass from a
varying fine-structure constant by sign structure, and the hydrogen-to-deuterium hyperfine ratio is
an internal control on universality. We treat $\varepsilon$ as a free amplitude to be fitted, not
predicted, and commit to no mechanism. The construction is a theoretical template for multi-band
programmes, not a new observational bound or an end-to-end survey forecast. A single band
deviating from its assigned weight falsifies the pattern.
\end{abstract}

\maketitle

% ---------------------------------------------------------------------------
\section{Introduction}\label{sec:intro}

Whether the fundamental constants vary in space or time has a long observational
literature~\cite{Uzan2003,Uzan2011}. Limits come from quasar absorption
spectra~\cite{Webb1999,MurphyWebbFlambaum2003}, from the cosmic microwave
background~\cite{PlanckVaryingConstants,HartChluba2018}, and from big-bang nucleosynthesis, the
Oklo natural reactor, and laboratory clock comparisons. Most of this work has concentrated on the
fine-structure constant; the electron mass has attracted more recent attention as well, including
in cosmological settings where a larger $m_e$ near recombination can alter the sound
horizon~\cite{Sekiguchi2021,HartChluba2018}. Those motivations are background. This paper does not
predict $\varepsilon$, does not commit to a mechanism for it, and does not claim to resolve any
cosmological tension.

The tightest existing limits on a fractional shift in $\mu \equiv m_p/m_e$ already come from
molecular spectroscopy, not from the radio continuum or from $21\,$cm absorption. Methanol
absorption toward PKS1830$-$211 bounds $|\Delta\mu/\mu|$ at a few$\times 10^{-7}$
\cite{KanekarMethanol2015,BagdonaiteMethanol2013} --- some thirty-five times tighter than the
best published $21\,$cm comparison reinterpreted as a bound on $\varepsilon$ below. That amplitude
superiority is settled; it is not what is at issue here. A single molecular comparison returns one
number against one free amplitude: every hypothesis with a free $\varepsilon$ accommodates it
equally well. What a multi-band radio pattern can supply, and a single line cannot, is a
\emph{falsifiable correlation}: fixed relative shifts across independent systematics, with a sign
structure that separates a varying $m_e$ from a varying $\alpha$.

Radio propagation has been used for varying constants before. Dispersion measures of localised fast
radio bursts have been used to bound variation in $\alpha$ and in the proton-to-electron mass
ratio~\cite{Kalita2024}, and the clustering of burst dispersion measures likewise~\cite{WangXia2025}.
Those analyses treat a single observable at a time. The construction used below, in which a
quantity reconstructed with laboratory constants is read against something that does not share its
dependence, is not new and is not confined to the radio bands. The thermal Sunyaev--Zel'dovich
effect carries the Thomson cross-section, $\sigma_T \propto \alpha^{2} m_e^{-2}$, so a
Comptonization inferred from it inherits a weight in exactly the sense of Eq.~\eqref{eq:weight}.
Compared against an independent X-ray determination of the same intracluster gas, that weight has
been used to bound variation in $\alpha$ over 119 clusters~\cite{LiuAlphaClusters2021}, and over
618 for spatial rather than temporal variation~\cite{deMartinoAlphaClusters2016}. What follows
applies the same logic to a different set of observables, chosen so that the weights differ in sign
as well as magnitude.

Different radio observables depend on the electron mass through different powers, so a universal
shift in $m_e$ does not move them together: it moves them by fixed, unequal, calculable amounts.
Section~\ref{sec:weights} derives five such weights as a statement of how a shift would
\emph{propagate}. Only two of those rows --- the $21\,$cm hyperfine frequency and Faraday rotation
--- are presently usable in a Gaussian multi-band analysis of the form
Eq.~\eqref{eq:ml}; the recombination-line, synchrotron and dispersion-measure rows fail that
presumption for reasons of kind rather than precision (Sec.~\ref{sec:sens}). The five-row vector is
therefore not a claim of five simultaneous measurements available today, and this paper is a
theoretical template for multi-band programmes, not a new observational bound or an end-to-end
survey forecast.

The argument is arranged as follows. Section~\ref{sec:setup} fixes notation and separates the two
kinds of observable involved. Section~\ref{sec:weights} derives the five weights.
Section~\ref{sec:alpha} shows how the pattern separates a varying electron mass from a varying
fine-structure constant, and examines an isotopic control that is easy to misread.
Section~\ref{sec:assume} states what the argument assumes. Section~\ref{sec:sens} works out which
rows can be entered today, the formal sensitivity of the usable pair, and where the amplitude
already stands from methanol.
% ---------------------------------------------------------------------------

\section{Setup}\label{sec:setup}

We consider a universal, species-independent fractional shift in the electron mass,
%
\begin{equation}
  m_e \;\longrightarrow\; m_e\,(1+\varepsilon), \qquad |\varepsilon| \ll 1 ,
  \label{eq:shift}
\end{equation}
%
holding all other constants fixed. We make no assumption about what produces
Eq.~\eqref{eq:shift} and we do not predict $\varepsilon$; it is a single free amplitude. The
content of this paper is the set of \emph{relative} responses, which do not depend on it.

\subsection{What is held fixed, and why the statement is dimensionless}

A shift in a dimensionful constant is not by itself physically meaningful: the magnitude and even
the sign of an inferred effect can depend on the choice of units, an objection raised explicitly
against earlier $21\,$cm work on varying constants~\cite{FlambaumPorsev2010,KhatriWandelt2007,
KhatriWandelt2010}. Equation~\eqref{eq:shift} must therefore be read with its reference declared.

We hold the fine-structure constant and the strong scale fixed, so that Eq.~\eqref{eq:shift} is a
statement about the dimensionless mass ratio
%
\begin{equation}
  \mu \equiv m_p/m_e, \qquad \delta\mu/\mu = -\varepsilon .
  \label{eq:mu}
\end{equation}
%
Every row below then restates without reference to a unit system. The two line frequencies become
dimensionless when referred to the proton mass, $\nu_{\rm hf}/m_p \propto \alpha^{4}\mu^{-2}$ and
$\nu_{\rm RRL}/m_p \propto \alpha^{2}\mu^{-1}$, and are observed as apparent redshift offsets, which
are dimensionless in any case. The dispersion and rotation rows are ratios of a reconstructed
quantity to the true one, with the same units above and below. The synchrotron row is dimensionless
when referred to the proton cyclotron frequency, $\nu_c/\nu_{c,p} \propto \mu$.

Recast this way the weight vector is unchanged, so nothing in Sec.~\ref{sec:weights} depends on the
unit convention. We continue to write the weights against $m_e$ because it is the more familiar
form, with Eq.~\eqref{eq:mu} understood throughout.

For an observable $\mathcal{O}$ we write its weight as the logarithmic derivative
%
\begin{equation}
  w_{\mathcal{O}} \;\equiv\; \frac{\partial \ln \mathcal{O}}{\partial \ln m_e},
  \qquad\text{so}\qquad
  \frac{\delta \mathcal{O}}{\mathcal{O}} = w_{\mathcal{O}}\,\varepsilon
  \;+\; \mathcal{O}(\varepsilon^2).
  \label{eq:weight}
\end{equation}

Two classes of observable respond differently, and the distinction matters for what is measurable.

Line frequencies are one class. For a single line a shift in the rest frequency is absorbed into
an apparent redshift: an observer measuring $\nu_{\rm obs}$ and assuming the laboratory
$\nu_{\rm rest}$ infers $1+z_{\rm app} = (1+z)(1+w\varepsilon)$. One line alone therefore carries
no information about $\varepsilon$. Two lines of different weight do, since their inferred
redshifts differ by a known factor.

Dispersion and rotation measures are the second class. These are not frequencies but quantities
reconstructed from a measured delay or rotation angle using laboratory constants. A shift in $m_e$
along the line of sight biases the reconstruction even when the physical column is unchanged.

The pattern derived below spans both classes, so it is not degenerate with redshift.

\section{The five weights}\label{sec:weights}

\subsection{Hydrogen $21\,$cm hyperfine: $w = +2$}

The ground-state hyperfine splitting is the Fermi contact interaction between the electron and
nuclear moments~\cite{WolfeBrownRoberts1976,Tzanavaris2005}. Its scale is the Rydberg reduced by
the ratio of the nuclear to the Bohr
magneton and by one further factor $\alpha^2$,
%
\begin{equation}
  \nu_{\rm hf} \;\propto\; \alpha^{2}\,
  \underbrace{\frac{\mu_N}{\mu_B}}_{\textstyle \propto\, m_e/m_p}\;
  \underbrace{\vphantom{\frac{\mu_N}{\mu_B}} R_\infty c}_{\textstyle \propto\, \alpha^{2} m_e}
  \;\propto\; g_p\,\alpha^{4}\,\frac{m_e^{2}}{m_p},
  \label{eq:hf}
\end{equation}
%
giving $w = +2$.

\subsection{Radio recombination lines: $w = +1$}

Transitions between high principal quantum numbers are pure Rydberg~\cite{GordonSorochenko2009},
$\nu = R_\infty c\,(n_1^{-2}-n_2^{-2})$ with $R_\infty \propto \alpha^{2} m_e$, so
%
\begin{equation}
  \nu_{\rm RRL} \;\propto\; m_e, \qquad w = +1 .
\end{equation}

\subsection{Dispersion measure: $w = -1$}

The group delay of a signal traversing a cold plasma is~\cite{LorimerKramer2004}
%
\begin{equation}
  t_{\rm DM} \;=\; \frac{e^{2}}{2\pi m_e c}\,\frac{1}{\nu^{2}} \int n_e \, dl .
  \label{eq:dm}
\end{equation}
%
The measured quantity is $t_{\rm DM}$; the reported dispersion measure is obtained by dividing
out~\cite{Kulkarni2020}
the \emph{laboratory} constant $e^{2}/2\pi m_e^{\rm lab}c$. If the true electron mass along the path
is $m_e^{\rm lab}(1+\varepsilon)$, the delay produced by a fixed physical column
$\int n_e\,dl$ is smaller by $(1+\varepsilon)^{-1}$, and the inferred column inherits that factor:
%
\begin{equation}
  \mathrm{DM}_{\rm inferred} = \left(\int n_e\,dl\right)(1-\varepsilon) + \mathcal{O}(\varepsilon^2),
  \qquad w = -1 .
\end{equation}
%
The bias is in the \emph{reconstruction}, not in the plasma. This row therefore tests the shift
against any independent determination of the same electron column.

\subsection{Synchrotron characteristic frequency: convention-dependent ($-1$ or $-3$)}

The characteristic frequency of synchrotron emission is~\cite{RybickiLightman1979}
%
\begin{equation}
  \nu_c \;=\; \frac{3}{2}\,\gamma^{2}\,\frac{eB}{2\pi m_e c}.
  \label{eq:sync}
\end{equation}
%
This is the one convention-dependent entry, and the convention is not a bookkeeping detail: it
changes the weight by a factor of three and must sit front and centre before the row is used.
Holding the magnetic field and the emitting particle's Lorentz factor fixed gives $w = -1$ from the
explicit $m_e$ in Eq.~\eqref{eq:sync}. Holding the particle's energy fixed instead lets
$\gamma = E/m_e c^{2}$ float, contributing a further $-2$, so that
%
\begin{equation}
  w_{\rm sync} = \begin{cases}
    -1, & \text{fixed } \gamma \text{ (fixed-field reading)},\\[2pt]
    -3, & \text{fixed } E \text{ (fixed-energy reading)}.
  \end{cases}
\end{equation}
%
Neither value is default physics. Real source modelling often labels the emitting population by
energy (injection spectra, cooling breaks, equipartition arguments), and under that labelling the
row is $-3$; under a fixed-$\gamma$ labelling it is $-1$. The choice must be fixed by the physical
situation before the row is read, and if the population cannot be labelled unambiguously the row is
ill-posed as a test of $\varepsilon$. For bookkeeping we quote $-1$ in Eq.~\eqref{eq:pattern} and
Table~\ref{tab:weights} under an explicit fixed-$\gamma$ convention; that is a declaration, not a
claim that fixed-$\gamma$ is the correct description of any particular source. We return to this in
Sec.~\ref{sec:assume}.

\subsection{Faraday rotation: $w = -2$}

The rotation measure of a magnetised plasma is~\cite{BrentjensDeBruyn2005}
%
\begin{equation}
  \mathrm{RM} \;=\; \frac{e^{3}}{2\pi m_e^{2} c^{4}} \int n_e B_\parallel \, dl ,
  \qquad w = -2 ,
\end{equation}
%
again a bias in the reconstructed quantity rather than in the medium.

\subsection{The pattern}

Collecting Eqs.~\eqref{eq:hf}--\eqref{eq:sync}:
%
\begin{equation}
  \boxed{\;
  \begin{aligned}
    w_{21} : w_{\rm RRL} &: w_{\rm DM} : w_{\rm sync} : w_{\rm RM} \\[2pt]
    &=\; +2 : +1 : -1 : -1 : -2
  \end{aligned}
  \;}
  \label{eq:pattern}
\end{equation}
%
Every entry follows from a textbook scaling once its convention is fixed; the synchrotron entry is
the sole convention-dependent weight, quoted here as $-1$ under the fixed-$\gamma$ declaration of
Sec.~\ref{sec:weights}. The amplitude $\varepsilon$ is common to all five and cancels from every
ratio, so Eq.~\eqref{eq:pattern} is a prediction about \emph{correlations} --- how a shift would
propagate across bands --- that can be tested without knowing the size of the shift.

\section{Separation from a varying fine-structure constant}\label{sec:alpha}

The competing hypothesis is a shift in $\alpha$ rather than $m_e$. The two are separable because
they produce different weight vectors across the same five bands. Working in Gaussian units, where
$e^{2} = \alpha \hbar c$ and hence $e \propto \alpha^{1/2}$, the same derivations give the
$\alpha$-weights in Table~\ref{tab:weights}.

\begin{table}[t]
\caption{Response of each observable to a fractional shift in the electron mass and in the
fine-structure constant, $w \equiv \partial \ln \mathcal{O} / \partial \ln X$. The two columns are
not proportional, so the hypotheses are distinguishable; the sign structure separates them without
reference to precision.}
\label{tab:weights}
\begin{ruledtabular}
\begin{tabular}{lcc}
observable & $w$ for $m_e$ & $w$ for $\alpha$ \\
\hline
$21\,$cm hyperfine          & $+2$ & $+4$   \\
radio recombination line    & $+1$ & $+2$   \\
dispersion measure (inferred) & $-1$ & $+1$   \\
synchrotron $\nu_c$ (fixed $\gamma$; $-3$ at fixed $E$) & $-1$ & $+1/2$ \\
Faraday rotation measure    & $-2$ & $+3/2$ \\
\end{tabular}
\end{ruledtabular}
\end{table}

The distinction is qualitative before it is quantitative. Every $\alpha$-weight is positive: a
shift in the fine-structure constant moves all five observables in the same direction. The
$m_e$-weights change sign across the set, raising two observables and lowering three. Measuring any
two bands of opposite $m_e$-weight therefore separates the hypotheses by the sign of the
correlation alone.

The reason is visible in the two reconstructed quantities. The dispersion delay of
Eq.~\eqref{eq:dm} carries $e^{2}/m_e \propto \alpha/m_e$, so $\alpha$ and $m_e$ enter with opposite
signs; the same holds for the rotation measure, where $e^{3}/m_e^{2} \propto \alpha^{3/2}/m_e^{2}$.
The line frequencies, by contrast, carry both constants with positive powers. No overall rescaling
maps one column of Table~\ref{tab:weights} onto the other: the ratios $w_\alpha / w_{m_e}$ across
the five rows are $2$, $2$, $-1$, $-1/2$ and $-3/4$.

\subsection{The hydrogen--deuterium control}

The hyperfine frequencies of hydrogen and deuterium both take the form of Eq.~\eqref{eq:hf} with
their own nuclear factors. Their ratio is therefore
%
\begin{equation}
  \frac{\nu_{\rm hf}({\rm H})}{\nu_{\rm hf}({\rm D})}
  \;=\; \frac{C_{\rm H}\, g_p / m_p}{C_{\rm D}\, g_d / m_d},
\end{equation}
%
in which both $\alpha$ and $m_e$ cancel identically, the angular-momentum coefficients $C$
depending only on the nuclear spins. Logarithmic differentiation confirms
$\partial \ln(\text{ratio})/\partial \ln m_e = \partial \ln(\text{ratio})/\partial \ln \alpha = 0$.

The ratio is therefore not a discriminant between a varying $m_e$ and a varying $\alpha$, since
both preserve it. What it tests is universality: whether whatever shifts does so identically for
the two isotopes. It excludes species-dependent variation, which is a useful null in its own right,
and it is easily misread as evidence for one constant over the other.

\section{What is assumed}\label{sec:assume}

The weights in Eq.~\eqref{eq:pattern} follow from standard atomic and plasma physics, but the test
built on them carries assumptions that should be stated.

\emph{Universality.} The shift is taken to act identically on every electron. A species- or
environment-dependent variation would break the pattern, and the hydrogen--deuterium ratio of
Sec.~\ref{sec:alpha} is the control for this.

\emph{Other constants fixed.} We vary $m_e$ alone. A simultaneous variation in $\alpha$ or in the
proton mass adds its own weights to each row, and the observed pattern would then be a
superposition. Because the weight vectors differ between $m_e$ and $\alpha$, a sufficiently
well-measured set of bands can in principle separate the components; with few bands it cannot.

\emph{Linearity.} Existing bounds on a varying electron mass from big-bang nucleosynthesis and
the microwave background place $|\varepsilon|$ at the percent level or
below~\cite{SetoToda2023,HartChluba2020}. We work to first order in $\varepsilon$. For
$|\varepsilon| \lesssim 10^{-2}$ the quadratic terms enter at the $10^{-4}$ level, below the
precision of any of the measurements considered here.

\emph{Constancy along the path.} Each observable integrates along a line of sight or over an epoch.
The weights apply cleanly if $\varepsilon$ is constant over the relevant interval. If it varies,
each row picks up its own weighting of the variation and the ratios are modified in a
row-dependent way. Comparisons at matched redshift are therefore preferable to comparisons across
widely separated epochs.

\emph{The synchrotron convention.} As noted in Sec.~\ref{sec:weights}, this row reads $-1$ or $-3$
according to whether the emitting population is labelled by Lorentz factor or by energy. The
$-1$ entry in Eq.~\eqref{eq:pattern} is a declared fixed-$\gamma$ convention, not a preferred
physical default. Real source modelling may make the labelling ambiguous, in which case the row
is ill-posed and should be dropped from any fit rather than entered under an unspoken choice.

\section{Sensitivity}\label{sec:sens}

Suppose each band $i$ yields a measured fractional shift $d_i$ with Gaussian uncertainty
$\sigma_i$, against the prediction $d_i = w_i \varepsilon$. The maximum-likelihood amplitude and
its error are
%
\begin{equation}
  \hat{\varepsilon} = \frac{\sum_i w_i d_i / \sigma_i^{2}}{\sum_i w_i^{2}/\sigma_i^{2}},
  \qquad
  \sigma_{\varepsilon} = \left( \sum_i \frac{w_i^{2}}{\sigma_i^{2}} \right)^{-1/2}.
  \label{eq:ml}
\end{equation}
%
The weights enter quadratically, so the high-$|w|$ rows dominate. Only two rows are presently
usable in Eq.~\eqref{eq:ml}, for reasons of kind set out immediately below; among those, the
$21\,$cm and Faraday rows are the extremes of the weight set and give
$\sigma_{\varepsilon} = \sigma/\sqrt{8} \simeq 0.35\,\sigma$ at common fractional precision
$\sigma$. That pair, not a five-band programme, is the realistic near-term test. For orientation
only, if all five bands were somehow measured at common precision one would have
$\sum_i w_i^{2} = 11$ and $\sigma_{\varepsilon} = \sigma/\sqrt{11} \simeq 0.30\,\sigma$; that figure
is arithmetic on the weight vector, not a survey forecast, and we do not treat it as a reachable
target.

Equation~\eqref{eq:ml} presumes each band supplies an independent Gaussian $\sigma_i$; three of the
five rows do not. Radio recombination frequencies are \emph{computed} from the Rydberg formula and
therefore already carry $m_e$, so an observed line compared against its own computed rest frequency
returns a velocity rather than a constraint; that row enters only in differential comparison
against a transition of different $m_e$ dependence. Synchrotron carries a formal degeneracy between
the field strength and the emitting population, broken in practice by assuming energy equipartition,
so its uncertainty is dominated by a modelling choice that is itself degenerate with the quantity
being weighed --- the resulting normalisation correction does not have a fixed sign, and the
$-1$ versus $-3$ convention of Sec.~\ref{sec:weights} may leave the row ill-posed for a real
source. The dispersion-measure row fails the same presumption, for the reason set out below: a
constant $\varepsilon$ is absorbed exactly by the fitted dispersion measure, so the row reports the
precision of an independent electron-column determination rather than a $\sigma_i$ of its own.

We therefore treat only the two rows referred to a fixed laboratory rest frequency ($21\,$cm and
Faraday rotation) as presently measurable. The five-row pattern of Sec.~\ref{sec:weights} is
unaffected: it is a statement about how a shift propagates, not a claim of five simultaneous
Gaussian measurements available today. What is reduced is the number of rows that can presently be
turned into a fit.

Pairs of opposite $m_e$-weight are preferable for a reason independent of their statistical
weight. A systematic common to both bands, entering as a shared multiplicative offset, contributes
equally to $d_{21}$ and $d_{\rm RM}$ and therefore cancels in the difference that carries the
signal, whereas it would be absorbed into $\hat{\varepsilon}$ by a same-sign pair. The sign
structure of Table~\ref{tab:weights} is thus useful twice: it separates the $m_e$ and $\alpha$
hypotheses, and it suppresses common-mode systematics.

\subsection{Where the rows currently stand}

The two presently measurable rows are quoted below with their statistical figures and, separately,
the systematic that actually binds each; in both cases the systematic exceeds the statistical
precision by roughly half an order of magnitude, and in each it arises from a different physical
effect. The dispersion-measure row is set out between them, not as a third constraint but because
its exclusion rests on a degeneracy the reader should be able to check rather than on an assertion.

The $21\,$cm row is constrained by comparing hyperfine absorption against optical metal lines in
the same absorber. The tightest such analysis reports $\Delta x / x = -(0.1 \pm 1.3) \times
10^{-6}$ over four systems at $1.17 < z < 1.56$, where $x \equiv g_p \alpha^{2}/\mu$
\citep{Rahmani2012}. Since $\mu = m_p/m_e$, holding $\alpha$ and the proton moment fixed gives
$\delta \ln x = -\delta \ln \mu = \varepsilon$, so that measurement is a measurement of the
amplitude itself:
%
\begin{equation}
  \varepsilon = -(0.1 \pm 1.3) \times 10^{-6} \quad (1.17 < z < 1.56),
  \label{eq:eps21}
\end{equation}
%
statistical only. An earlier analysis of the same combination
quotes a maximum systematic of $6.7 \times 10^{-6}$ \citep{Kanekar2010}, arising because the
radio and optical absorption need not sample the same gas along a finite-width sightline --- which
is why the tighter analysis restricts itself to milliarcsecond-unresolved background sources.
Added in quadrature the two give $\varepsilon = -(0.1 \pm 6.8) \times 10^{-6}$, or
$|\varepsilon| < 1.4 \times 10^{-5}$ at $2\sigma$. This is a reinterpretation of a published
measurement under the assumptions of Sec.~\ref{sec:assume}, not a new one. The systematic exceeds
the statistical error by a factor of five --- exactly the situation a ratio test across bands with
independent systematics is meant to improve on.

The dispersion-measure row has by far the best raw precision and, for the same reason as the
recombination-line row, the least of it is usable. Median DM uncertainties reach
$\sim\!10^{-5}\,\mathrm{pc\,cm^{-3}}$ in low-frequency monitoring, which against typical
pulsar-timing-array dispersion measures of $10$--$100\,\mathrm{pc\,cm^{-3}}$ is a fractional
precision of $10^{-7}$--$10^{-6}$. That precision does not transfer to $\varepsilon$. A pulsar
timing model fits an infinite-frequency arrival time and a dispersion measure from
$t(\nu) = t_{\infty} + K\,\mathrm{DM}/\nu^{2}$, and a universal shift $m_e \to m_e(1+\varepsilon)$
rescales the delay's coefficient while leaving its $\nu^{-2}$ shape untouched. The fit therefore
absorbs it exactly, returning $\mathrm{DM}_{\rm fit} = (\int n_e\,dl)/(1+\varepsilon)$ with
$t_{\infty}$ unchanged and no timing residual at any frequency coverage. A constant $\varepsilon$
is thus degenerate with the fitted dispersion measure, and improving the DM measurement improves
$\mathrm{DM}_{\rm fit}$ while saying nothing about the shift. The row constrains $\varepsilon$ only
in the sense already stated in Sec.~\ref{sec:weights}: against an independent determination of the
same electron column, whose precision, not the DM precision, sets $\sigma_\varepsilon$.

That route has been taken. The fast-radio-burst analyses cited in Sec.~\ref{sec:intro} compare the
observed dispersion measure against the value predicted from $\Omega_b h^{2}$ and an assumed
intergalactic baryon fraction, which is exactly the external column referred to above. Seventeen
bursts with supernovae in a runaway-dilaton parametrization give $\Delta\alpha/\alpha$ at the
$10^{-2}$ level~\cite{Lemos2025}; 50 bursts over $0.004 < z < 1.02$ tighten this to
$\Delta\alpha/\alpha \simeq 2 \times 10^{-5}$~\cite{Kalita2024}. One caution about reading those
results as mass constraints. The quoted $\Delta\mu/\mu \simeq -1 \times 10^{-5}$ is not independent
of the $\alpha$ result but derived from it, through an assumed relation between mass and coupling
variation, so it does not constrain $\varepsilon$ at fixed $\alpha$ --- the case considered here.
The corresponding mass-only analysis has not to our knowledge been performed, and its limiting
uncertainty would be the baryon fraction rather than the dispersion measures.

What the timing literature does bound is the \emph{variation} of $\varepsilon$. Mis-estimated
dispersion produces infinite-frequency arrival-time offsets of order $20\,\mu\mathrm{s}$ for
high-DM pulsars, correlated on month timescales
\citep{JonesNANOGrav2017,NANOGravDM2023}; read as a bound on epoch-to-epoch drift this gives
$\sim\!2\times10^{-6}$ at high DM and low frequency, falling to $\sim\!3\times10^{-4}$ at
$1400\,\mathrm{MHz}$ and low column. We therefore do not fold the $10^{-7}$ figure into any
sensitivity estimate, and we do not treat this row as carrying a constant-$\varepsilon$ constraint
at all.

The Faraday row is set by rotation-measure catalogue precision. The deepest current grid reports a
median RM uncertainty of $0.06\,\mathrm{rad\,m^{-2}}$ across $2461$ extragalactic sources, with a
possible residual systematic of up to $0.3\,\mathrm{rad\,m^{-2}}$ after ionospheric correction
\citep{OSullivanLoTSS2023}.

\subsection{Where the amplitude actually stands}

The strongest existing bound on $\varepsilon$ comes from none of these rows, and it is stated first
so that the paper is not read as competing on amplitude. Methanol absorption in the $z = 0.88582$
lens toward PKS1830$-$211 gives $|\Delta\mu/\mu| \lesssim 4 \times 10^{-7}$ at $2\sigma$ over
$0 < z \leq 0.886$ \citep{KanekarMethanol2015}. The methanol transitions carry $\mu$ alone --- no
$\alpha$, no proton moment --- so that limit constrains $\varepsilon$ under weaker assumptions than
Eq.~\eqref{eq:eps21} requires, and it is some thirty-five times tighter. An independent analysis of
the same source at three telescopes, using ten transitions, agrees:
$\Delta\mu/\mu = (-1.0 \pm 0.8_{\rm stat} \pm 1.0_{\rm sys}) \times 10^{-7}$
\citep{BagdonaiteMethanol2013}, which is $|\Delta\mu/\mu| < 3.6 \times 10^{-7}$ at $2\sigma$ and is
likewise dominated by its systematic. Matching that figure with the $21\,$cm and Faraday pair would
require a fractional precision of $\sigma \simeq 1.1 \times 10^{-6}$ per band; the two rows that can
presently be entered do not improve the limit on the amplitude, and we do not claim that they do.

What they supply instead is something a single molecular comparison cannot. One transition returns
one number against one free amplitude $\varepsilon$; the fit is exact and nothing is left over, so
every hypothesis with a free amplitude accommodates the result equally well. Two rows of different
weight return two numbers for the same amplitude, leaving one residual degree of freedom that the
hypotheses do not share --- and it is that residual, not the amplitude, that Eq.~\eqref{eq:ml}
calculates the sensitivity of. Because the count does not involve $\sigma$, a tighter bound on the
amplitude does not diminish the pattern test. The methanol analyses already illustrate the same
point that limits the $21\,$cm row: a tighter statistical limit from one line pair was set aside in
favour of the weaker figure from three transitions whose profiles agree, on the grounds that only
those demonstrably sample the same gas. Both of the best available constraints on $\varepsilon$ are
limited by sightline structure rather than by photon noise.

This paper is a theoretical template for multi-band programmes: it supplies weights, a falsification
rule, and an honest account of which rows enter a Gaussian fit today. It is not a new observational
bound, and it is not an end-to-end survey forecast.

\section{Falsification}

The pattern is falsified by any single band deviating from its assigned weight, once two or more
usable bands are measured along the same sightline or at matched redshift. With only the $21\,$cm
and Faraday rows presently enterable, the near-term falsifier is a pair whose residual is
inconsistent with a common $\varepsilon$; the full five-weight vector remains the theoretical
target if and when the other rows become well posed.

Two outcomes distinguish the hypothesis from the alternatives. A pattern consistent with
$+2:+1:-1:-1:-2$ across the bands that can be entered supports a universal shift in $m_e$. A
pattern with the same signs but different magnitudes is evidence for a mixed variation, or for a
variation in $\alpha$, and the weight vectors separate the cases. Uncorrelated shifts, or a single
anomalous band with the others null, are not evidence for this hypothesis and should be treated as
instrumental or astrophysical in origin.

\section{Conclusion}

A universal fractional shift in the electron mass would produce a correlated signature across radio
observables whose relative sizes are fixed by atomic physics. The ratios
$+2:+1:-1:-1:-2$ state how such a shift would propagate; they are independent of the amplitude, so
the structure can be tested without a prediction for that amplitude and without commitment to any
mechanism producing it. The pattern differs from that produced by a varying fine-structure
constant, and the hydrogen--deuterium hyperfine ratio provides an internal control on universality
that both hypotheses pass and a species-dependent variation fails.

Only two of the five rows --- $21\,$cm and Faraday rotation --- can be entered in a Gaussian
multi-band fit at present; the other three remain theoretical propagation channels. Methanol
absorption already bounds the amplitude some thirty-five times more tightly than the best published
$21\,$cm comparison ($|\varepsilon| < 1.4 \times 10^{-5}$ over $1.17 < z < 1.56$, systematic five
times the statistical error). Neither figure is improved here, and this paper is not a survey
forecast. The case for the pattern is not that it measures $\varepsilon$ better, but that a limit
on $\mu$ is a single number, which cannot say whether the shift lies in the electron mass or in the
coupling, nor whether every band shares it. Fixed weights of two signs answer both questions; a
single band off its assigned weight falsifies the pattern.

\bibliographystyle{apsrev4-2}
\bibliography{refs}

\end{document}
